\documentclass[11pt]{article}
\usepackage[margin=1in]{geometry}
\usepackage{amsmath,amsthm,amssymb}
\usepackage{hyperref,enumitem}

\newtheorem{theorem}{Theorem}[section]
\newtheorem{lemma}[theorem]{Lemma}
\newtheorem{corollary}[theorem]{Corollary}
\newtheorem{definition}[theorem]{Definition}
\newtheorem{remark}[theorem]{Remark}
\newtheorem*{openproblem}{Open Problem}

\title{\textbf{A Ring Lemma for Band-Limited Vorticity Direction\\
and a Conditional Spectral Non-Dispersal Criterion}\\[6pt]
\large on the Three-Dimensional Torus}
\author{Jonathan Robert Simons\\
Savannah, Georgia, USA\\
\texttt{simonsmedicalinnovations@gmail.com}}
\date{Preprint \quad|\quad Submission draft}

\begin{document}
\maketitle

\begin{center}
\fbox{\parbox{0.92\linewidth}{\centering\small
\textbf{Disclaimer.} The main regularity statement is \emph{conditional}
on Spectral Non-Dispersal (SND). Unconditional SND for arbitrary large $H^1$
data remains open.}}
\end{center}
\vspace{0.75em}

\begin{abstract}
We record two fluids-only tools for the incompressible Navier--Stokes equations on
$\mathbb{T}^3$. First, a \emph{Ring Lemma}: if the velocity is Fourier-supported in a
single Littlewood--Paley shell $S_{j^*}$, then on the set where vorticity is comparable
to its $L^2$ norm the vorticity direction satisfies
$\|\nabla\xi_0\|_{L^\infty(E_c)}\le C\,2^{j^*}$. Second, a \emph{Spectral Non-Dispersal}
(SND) condition requiring that a positive fraction of enstrophy remain in a dominant
shell, $\inf_t J(t)/X(t)\ge c_*>0$. Under SND we outline a conditional route to
$H^1$ control via the Ring Lemma in the concentrated regime and enhanced dissipation
in the spread regime. We do not claim unconditional global regularity.
\end{abstract}

\section{Setup}

Let $u$ be a Leray--Hopf solution of
\[
\partial_t u+(u\cdot\nabla)u+\nabla p=\nu\Delta u,\qquad \operatorname{div}u=0
\]
on $\mathbb{T}^3$, with $u_0\in H^1$ divergence-free. Write $\Delta_j$ for dyadic
Littlewood--Paley projections and
\begin{align*}
X_j(t)&=2^{2j}\|\Delta_j u(t)\|_{L^2}^2,\qquad
X(t)=\sum_j X_j(t)=\|\nabla u(t)\|_{L^2}^2,\\
J(t)&=\max_j X_j(t),\qquad
\rho(t)=J(t)/X(t),\qquad
j^*(t)=\operatorname{argmax}_j X_j(t).
\end{align*}

\begin{definition}[SND]
A solution satisfies \emph{Spectral Non-Dispersal} on an interval $I$ if there exists
$c_*>0$ such that
\[
\inf_{t\in I}\frac{J(t)}{X(t)}\ge c_*.
\]
Equivalently, the dominant shell always carries at least fraction $c_*$ of the enstrophy.
\end{definition}

\begin{remark}
Some companion notes use an upper-bound convention $\sup_j a_j\le\rho_0$ for normalized
shell energies $a_j$. Those are related but not identical; this note freezes the
lower-bound form above.
\end{remark}

\section{The Ring Lemma}

\begin{lemma}[Ring Lemma]
Suppose $u\in H^1(\mathbb{T}^3)$ has Fourier support in a single shell
$S_{j^*}=\{k\in\mathbb{Z}^3:|k|\sim 2^{j^*}\}$. Let $\omega=\operatorname{curl}u$,
$\xi_0=\omega/|\omega|$ where $\omega\neq 0$, and
\[
E_c=\{x:|\omega(x)|\ge c\|\omega\|_{L^2}\}
\]
for a fixed $c\in(0,1)$. Then $|E_c|>0$ for $c$ small enough, and
\[
\|\nabla\xi_0\|_{L^\infty(E_c)}\le C\,2^{j^*}
\]
with $C$ depending only on $c$ and Bernstein constants on $\mathbb{T}^3$.
The bound persists while Fourier support remains in $S_{j^*}$.
\end{lemma}

\begin{proof}[Proof sketch]
On a single shell, Bernstein yields
$\|\nabla\omega\|_{L^\infty}\lesssim 2^{j^*}\|\omega\|_{L^\infty}
\lesssim 2^{2j^*}\|\omega\|_{L^2}$ (up to absolute constants).
On $E_c$,
\[
|\nabla\xi_0|\le\frac{|\nabla\omega|}{|\omega|}
+\frac{|\nabla|\omega||}{|\omega|}
\le\frac{2\|\nabla\omega\|_{L^\infty}}{c\|\omega\|_{L^2}}
\lesssim\frac{2^{j^*}}{c}.
\]
Measure of $E_c$ follows by Chebyshev. Persistence follows from persistence of
Fourier support under the (smooth) dynamics on short times, or under an augmented
parabolic approximation.
\end{proof}

\begin{corollary}[Band-limited CF control]
Under the Ring Lemma hypotheses, $\nabla\xi_0$ is bounded on $E_c$ by a multiple of the
shell frequency. In that band-limited regime, geometric vorticity-direction criteria of
Constantin--Fefferman type apply with an explicit constant depending on $j^*$
\cite{CF1993}.
\end{corollary}

\section{Conditional regularity under SND}

\begin{theorem}[Conditional $H^1$ control]
Let $u_0\in H^1(\mathbb{T}^3)$ be divergence-free and let $u$ be an associated
Leray--Hopf solution. Suppose SND holds on $[0,T]$ with constant $c_*>0$. Then, in the
concentrated regime where a dominant shell persists,
\[
u\in L^\infty([0,T];H^1)\cap L^2([0,T];H^2)
\]
provided the Ring Lemma may be applied to a band-limited approximation of the dominant
shell contribution and standard dissipative estimates control the complementary spread
regime. In particular, no finite-time singularity occurs on $[0,T]$ under these
hypotheses.
\end{theorem}

\begin{proof}[Proof outline]
\textbf{Concentrated regime.} SND forces $J\ge c_*X$, so a dominant scale $j^*$ carries a
fixed enstrophy fraction. Approximating that contribution by a band-limited field and
applying the Ring Lemma yields a vorticity-direction bound of size $O(2^{j^*})$ on $E_c$,
feeding a CF-type geometric estimate that prevents vortex-stretching blowup while the
band persists.

\textbf{Spread regime.} When enstrophy is dispersed across many shells, shell-spread
Poincar\'e / enhanced dissipation estimates give additional $H^1$ damping
(companion quantitative notes). Finite occupation time of extreme spread regimes follows
from energy identities.

\textbf{Glue.} Alternation between regimes under SND keeps $\|\nabla u\|_{L^2}$ from
blowing up on $[0,T]$. Details of constants and approximation are deferred to the longer
conditional framework preprint; the present note isolates the geometric lemma and the
conditional claim.
\end{proof}

\begin{remark}[What is not claimed]
We do \emph{not} prove SND for arbitrary large data.
We do not identify SND with arithmetic operators $1/\gcd$ or even-sum
prime-partition conditions.
\end{remark}

\section{Open problem}

\begin{openproblem}
Prove or disprove that every Leray--Hopf solution with $u_0\in H^1(\mathbb{T}^3)$
satisfies SND on $[0,\infty)$ for some $c_*=c_*(u_0,\nu)>0$.
\end{openproblem}

\begin{thebibliography}{9}
\bibitem{CF1993}
P.~Constantin and C.~Fefferman,
\textit{Direction of vorticity and the problem of global regularity for the
Navier--Stokes equations},
Indiana Univ.\ Math.\ J.\ \textbf{42} (1993), 775--789.

\bibitem{BKM1984}
J.~T.~Beale, T.~Kato, and A.~Majda,
\textit{Remarks on the breakdown of smooth solutions for the 3-D Euler equations},
Comm.\ Math.\ Phys.\ \textbf{94} (1984), 61--66.
\end{thebibliography}

\end{document}
